\scanpage{27}
\mspage{12}
L'extension $\mathcal E=\mathcal E_{0,3}$ de $\Gamma=\Gamma_{0,3}=\underbrace{\unsure{\mathfrak S_{\{0,1,\infty\}}|\mathcal S_{\{0,1,\infty\}}}}_{\unsure{\mathfrak S_3|\mathcal S_3}}\times\mathbb Z/2\mathbb Z$
\margin{engendré par $\tau:\ z\mapsto\bar z$}
par
\[
\pi=\pi_{0,3}=\{l_0,l_1,l_\infty\mid l_\infty l_1l_0=1\}=\pi_1(\underbrace{\mathbb P^1_{\mathbb C}\setminus\{0,1,\infty\}}_{U=U_{0,3}},-\bar j).
\]
Les classes de conjugaison \deleted{\unsure{au dessus des}} de scindages partiels maximaux correspondent aux points fixes de $\dot\rho,\dot\sigma_0,\dot\sigma_1,\dot\sigma_\infty$ opérant sur $X_{0,3}$, et \inserted{(éventuelles)} aux comp. connexes des ens. d'invariants de $\deleted{\unsure{\dot\rho}}\tau,\dot\sigma_0\tau,\dot\sigma_1\tau,\dot\sigma_\infty\tau$. Les pts fixes de $\rho$ sont $P$ et $\bar P$ dont le stabilisateur est d'ordre $6$, \inserted{savoir $\varphi(\unsure{\mathfrak S_3|\mathcal S_3})$} formé de $1,\dot\rho,\dot\rho^2,\dot\sigma_1\tau,\dot\sigma_2\tau,\dot\sigma_3\tau$. On a étudié le relèvement correspondant \inserted{à $P$} (de $\varphi(\unsure{\mathfrak S_3|\mathcal S_3})$), on trouve des relèvements ($\pi$-conjugués \inserted{entre eux}) relatifs à $\bar P$, s'obtiennent en \deleted{\unsure{conj.}} prenant un élémt de $\mathcal E$ \deleted{\unsure{qui}} qui soit au dessus d'un élémt de $\Gamma$ envoyant $P$ en $\bar P$ (savoir $\dot\sigma_0,\dot\sigma_1,\dot\sigma_\infty,\tau,\tau\rho,\tau\rho^2$) et en conjuguant par cet élément. On peut donc prendre \deleted{$\unsure{\sigma_0,\tau}$} un des éléments $\sigma_0,\sigma_1,\sigma_\infty$ (ou $\varepsilon_0,\varepsilon_1,\varepsilon_\infty$) ou $\tau_0,\tau_1,\tau_\infty$ ou les produits des précédents par \deleted{\unsure{$\tau$}} $\rho$ ou $\rho^2$\ldots

\unsure{Le|Les} \inserted{seul} pt fixe de $\sigma_0$ est $2$, dont le stabilisateur est le groupe $\simeq\mathbb Z/2\times\mathbb Z/2$ engendré par $\sigma_0$ et $\tau$. On a un relèvement de ce sous-groupe
\mspage{26}
